\documentclass[12pt]{article}

\usepackage[margin=1in]{geometry}
\usepackage{amsmath,amssymb}
\usepackage{booktabs}
\usepackage{graphicx}
\usepackage{natbib}
\usepackage{setspace}
\usepackage{hyperref}
\usepackage[T1]{fontenc}
\usepackage{lmodern}
\usepackage{threeparttable}

\hypersetup{
    colorlinks=true,
    linkcolor=blue!60!black,
    citecolor=blue!60!black,
    urlcolor=blue!60!black,
}

\doublespacing

\title{Online Appendix for ``Why Retirees Don't Annuitize: A Unified Lifecycle Model''}
\author{Derek Tharp}
\date{\today}

\begin{document}

\maketitle

\appendix
\setcounter{section}{0}
\renewcommand{\thesection}{\Alph{section}}


%% ================================================================
%%  A. MODEL DETAILS
%% ================================================================
\section{Model Details}
\label{app:model}

\subsection{Health transition matrices}
\label{app:health_transitions}

Health transitions are calibrated to the RAND HRS longitudinal file (waves 4--16), with the five self-reported health categories mapped to three states: Good (Excellent/Very Good), Fair (Good), and Poor (Fair/Poor). Transition probabilities are age-dependent, estimated in five-year age bands from 65--69 through 95+.

The transition matrices exhibit two standard features. First, health deterioration is more likely than improvement: transition rates from Good to Fair exceed those from Fair to Good at all ages. Second, deterioration accelerates with age: the probability of remaining in Good health falls from approximately 0.80 at ages 65--69 to 0.55 at ages 90+.

\subsection{Medical expenditure calibration}

Log out-of-pocket medical expenditures follow:
\[
\ln m_t = \mu_m^{\text{base}} + g \cdot (t - 65) + \delta_\mu(H_t) + \big[\sigma_m^{\text{base}} + \delta_\sigma(H_t)\big] \cdot \varepsilon_t, \quad \varepsilon_t \sim N(0,1),
\]
where $\mu_m^{\text{base}} = 7.037$ is the log mean at age 65 for Fair health, $g = 0.0652$ is the annual growth rate in the log mean, $\delta_\mu(H)$ shifts the log mean by health state $[-0.5, 0, 0.7]$ for Good, Fair, and Poor respectively, $\sigma_m^{\text{base}} = 1.4$ is the base log standard deviation, and $\delta_\sigma(H)$ shifts the log standard deviation by health $[-0.2, 0, 0.2]$.

\begin{table}[htbp]
\centering
\caption{Medical Expenditure Moments: Model vs.\ Target}
\label{tab:medical_moments}
\begin{tabular}{lrr}
\toprule
Moment & Model & Target (Jones et al.\ 2018) \\
\midrule
Mean OOP, age 70, Fair health & \$4,201 & \$4,200 \\
Mean OOP, age 100, Fair health & \$29,703 & \$29,700 \\
95th pctile, age 100, Fair health & \$111,502 & \$111,200 \\
\bottomrule
\end{tabular}
\end{table}

The lognormal specification produces a right-skewed distribution consistent with the heavy tails observed in medical expenditure data. The log standard deviation of 1.4 implies a coefficient of variation of approximately 2.5, which matches the empirical distribution of out-of-pocket spending in the HRS and MEPS.

\subsection{Survival calibration}
\label{app:survival}

Baseline survival probabilities use the SSA administrative period life table employed by \citet{lockwood2012}. The cumulative death probabilities at ages 65, 70, 75, \ldots are taken directly from Lockwood's replication code (BAP\_sim2.m) and interpolated linearly at intermediate ages.

Health-adjusted survival follows the proportional hazard specification $s(t, H) = s_{\text{base}}(t)^{\mu(H)}$, where $\mu(H)$ is the health-specific hazard multiplier. At baseline, $\mu = [0.50, 1.0, 3.0]$. This implies that an individual in Good health at age 80---where $s_{\text{base}}(80) \approx 0.96$---survives with probability $0.96^{0.5} \approx 0.98$, while an individual in Poor health survives with probability $0.96^{3.0} \approx 0.89$.

The constant-multiplier specification is a simplification. The HRS data show that the health-mortality gradient compresses with age: the ratio of Poor-to-Fair mortality is 3.29 at ages 65--74, 2.77 at 75--84, and 1.82 at 85+. Age-varying multipliers would slightly attenuate the Reichling--Smetters mechanism at advanced ages. We use constant multipliers for tractability and because the annuitization decision is made at age 65, when the gradient is widest.


%% ================================================================
%%  B. GRID CONVERGENCE
%% ================================================================
\section{Grid Convergence}
\label{app:grid_convergence}

Table~\ref{tab:grid_convergence} reports predicted ownership under the full model for different grid resolutions. The three grids---wealth ($n_W$), annuity income ($n_A$), and annuitization fraction ($n_\alpha$)---have partially offsetting effects on ownership. Coarser annuity income grids smooth over the dip in value at intermediate $\alpha$, biasing ownership upward. Coarser $\alpha$ grids miss beneficial annuitization fractions, biasing ownership downward. Both converge as resolution increases.

\begin{table}[htbp]
\centering
\caption{Grid Convergence}
\label{tab:grid_convergence}
\begin{tabular}{lccc}
\toprule
Resolution & $n_W \times n_A \times n_\alpha$ & Ownership (\%) & $\Delta$ from production \\
\midrule
Coarse & $30 \times 10 \times 21$ & 22.1 & $+20.7$ pp \\
 & $40 \times 15 \times 31$ & 13.6 & $+12.2$ pp \\
Medium & $60 \times 20 \times 51$ & 6.5 & $+5.1$ pp \\
 & $60 \times 25 \times 101$ & 2.8 & $+1.4$ pp \\
\textbf{Production} & $\mathbf{80 \times 30 \times 101}$ & \textbf{1.4} & --- \\
Fine & $80 \times 30 \times 201$ & 1.6 & $+0.2$ pp \\
 & $100 \times 40 \times 101$ & 1.4 & $0.0$ pp \\
 & $100 \times 40 \times 201$ & 1.4 & $0.0$ pp \\
\bottomrule
\end{tabular}
\begin{tablenotes}
\small
\item All specifications use baseline parameters ($\gamma = 2.5$, DFJ bequests, MWR $= 0.82$, $\pi = 2\%$, hazard multipliers $[0.50, 1.0, 3.0]$).
\end{tablenotes}
\end{table}

The production grid (80 $\times$ 30 $\times$ 101) yields 1.4\% ownership. Doubling the wealth and annuity grids to 100 $\times$ 40 while holding $n_\alpha = 101$ produces the same 1.4\%. Doubling $n_\alpha$ to 201 also produces 1.4\%. Convergence within $\pm 0.2$ percentage points is achieved at the production resolution.


%% ================================================================
%%  C. BEQUEST SPECIFICATION DETAILS
%% ================================================================
\section{Bequest Specification}
\label{app:bequest}

The De Nardi--French--Jones (DFJ) bequest utility is $v(b) = \theta (b + \kappa)^{1-\gamma}/(1-\gamma)$. The parameter $\kappa$ controls the luxury-good nature of bequests. When $\kappa$ is large relative to wealth $b$, marginal bequest utility $v'(b) = \theta(b + \kappa)^{-\gamma}$ is nearly constant across wealth levels, making bequests a luxury good: the bequest-to-wealth ratio rises with wealth.

The baseline parameters ($\theta = 56.96$, $\kappa = \$272{,}628$) are from \citet{lockwood2012}, who estimated them from HRS data under the DFJ specification. At these values, an individual with \$50,000 in wealth has $v'(b) \approx \theta \cdot (322{,}628)^{-2.5}$, which is negligible compared to the marginal utility of consumption. The bequest motive binds meaningfully only for individuals with wealth approaching $\kappa$.

\paragraph{Homothetic bequest anomaly.} Combining the DFJ bequest intensity $\theta = 56.96$ with a homothetic specification ($\kappa = \$10$) produces anomalous results. At $\kappa = \$10$, marginal bequest utility near $b = 0$ is $\theta \cdot 10^{-2.5} = 56.96 \times 3.16 \times 10^{-3} \approx 0.18$, which is extreme relative to marginal consumption utility. This creates a narrow precautionary buffer---individuals avoid annuitizing not to protect bequests, but to avoid the near-zero-wealth states where $v'(b)$ diverges. The predicted ownership under this misspecification (60.8\%) is higher than the no-bequest case (15.7\%), because the buffer effect dominates the bequest effect.

This anomaly illustrates why $\theta$ and $\kappa$ must be interpreted as a pair, not independently. The DFJ estimates from \citet{lockwood2012} are jointly calibrated; using one with the other's value replaced is not a valid counterfactual.


%% ================================================================
%%  D. PASHCHENKO COMPARISON DETAILS
%% ================================================================
\section{Comparison with Pashchenko (2013)}
\label{app:pashchenko}

\citet{pashchenko2013} conducted a sequential decomposition with four channels: pre-existing annuitization (Social Security and DB pensions), bequest motives, minimum purchase requirements, and pricing loads. Her full model predicted approximately 20\% participation, roughly four times the observed rate of 5\%.

We replicate her channel set within our framework. Our replication differs from hers in three respects: we use DFJ luxury-good bequests rather than the homothetic specification she employed; we use a continuous annuitization fraction rather than a binary purchase decision; and we do not model housing illiquidity. Despite these differences, the qualitative pattern matches.

\begin{table}[htbp]
\centering
\caption{Pashchenko (2013) Channels in Our Framework}
\label{tab:pashchenko}
\begin{tabular}{lcc}
\toprule
Model Specification & Ownership (\%) & $\Delta$ \\
\midrule
Yaari benchmark (SS on) & 100.0 & --- \\
+ Bequest motives (DFJ) & 100.0 & +0.0 pp \\
+ Minimum purchase (25K) & 74.1 & -25.9 pp \\
+ Pricing loads (MWR=0.85) & 59.5 & -14.7 pp \\
\midrule
\textit{Channels omitted by Pashchenko:} & & \\
+ Medical costs + R-S correlation & 38.2 & -21.3 pp \\
+ Survival pessimism & 13.7 & -24.4 pp \\
+ Full loads + Inflation & 8.5 & -5.3 pp \\
\midrule
Observed (Lockwood 2012) & 3.6 & --- \\
\bottomrule
\end{tabular}
\begin{tablenotes}
\small
\item Steps 0--3 incorporate the channels Pashchenko (2013) identified
(SS, bequests, minimum purchase, pricing loads) into our unified model.
Steps 4--6 add channels not in her framework. Note: this is not a
replication of her model (different preferences, health dynamics,
solution method). Housing illiquidity is not modeled; see text.
\end{tablenotes}
\end{table}


Under Pashchenko's channels alone (Steps 0--3), our model predicts 66.6\% ownership. The higher prediction relative to Pashchenko's 20\% reflects differences in model specification, particularly the continuous annuitization margin. With a binary purchase decision and housing illiquidity (which Pashchenko includes and we do not), predicted participation would be lower.

The channels Pashchenko omitted---medical expenditure risk with health-mortality correlation (Step 4) and inflation erosion (Step 5)---reduce ownership from 66.6\% to 23.9\%. These two channels account for a larger ownership reduction (42.7 percentage points) than Pashchenko's four channels combined (19.5 percentage points from Steps 0 to 3).


%% ================================================================
%%  E. DIA/QLAC EXTENSION
%% ================================================================
\section{Deferred Income Annuity Extension}
\label{app:dia}

Deferred income annuities (DIAs) begin payments at a future age, typically 80 or 85. By concentrating payments in late life---when longevity risk is greatest---they offer higher per-dollar payouts. A DIA purchased at 65 with payments beginning at 80 pays roughly 2.8 times the per-period income of an immediate annuity (SPIA) per dollar of premium, and a DIA-85 pays roughly 5.9 times as much.

However, deferred products carry lower money's worth ratios. \citet{wettstein2021} estimate MWRs of 0.50 for DIA-80 and 0.45 for DIA-85, compared to 0.82 for SPIAs. The lower MWR reflects the longer deferral period (more discounting), smaller market (less competition), and higher sensitivity to mortality projection errors.

\begin{table}[htbp]
\centering
\caption{SPIA vs.\ Deferred Income Annuity (DIA) Comparison}
\label{tab:dia}
\begin{tabular}{lccc}
\toprule
 & SPIA & DIA-80 & DIA-85 \\
MWR & 0.87 & 0.50 & 0.45 \\
Payout rate & 0.0623 & 0.1659 & 0.3485 \\
\midrule
\multicolumn{4}{l}{\textit{No bequest}} \\
\quad Ownership (\%) & 1.6 & 0.0 & 0.4 \\
\quad CEV at \$200K (\%) & 0.00 & 0.00 & 0.00 \\
\quad Optimal $\alpha$ at \$200K & 0.0\% & 0.0\% & 0.0\% \\
\midrule
\multicolumn{4}{l}{\textit{Moderate (DFJ)}} \\
\quad Ownership (\%) & 0.0 & 0.0 & 0.0 \\
\quad CEV at \$200K (\%) & 0.00 & 0.00 & 0.00 \\
\quad Optimal $\alpha$ at \$200K & 0.0\% & 0.0\% & 0.0\% \\
\bottomrule
\end{tabular}
\begin{tablenotes}
\small
\item DIA MWR from Wettstein et al.\ (2021). All models include
medical costs, R-S correlation, pricing loads, and 2\% inflation.
CEV evaluated at good health (H=1).
\end{tablenotes}
\end{table}


Table~\ref{tab:dia} compares predicted ownership and welfare across product types. All models include medical costs, health-mortality correlation, pricing loads, and 2\% inflation.

Under no bequests, SPIA ownership is 20.6\%, DIA-80 is 8.4\%, and DIA-85 is 6.6\%. The lower DIA ownership reflects the poor MWR: the actuarial advantage of deferred payout is more than offset by the pricing disadvantage. CEV at \$200,000 wealth is zero for all three products.

With DFJ bequests, SPIA ownership falls to 5.9\% while DIA ownership is unchanged at 8.4\% (DIA-80) and 6.6\% (DIA-85). DIAs are less affected by bequests because the premium is spent at age 65 but generates income only in late life, creating a natural separation between bequeathable wealth and annuity purchase.

The DIA results suggest that product innovation alone is unlikely to resolve low annuity demand. The channels suppressing SPIA demand---valuation risk, inflation erosion, bequest motives---apply with similar force to deferred products, and the inferior pricing of DIAs partially offsets their actuarial advantage.


%% ================================================================
%%  F. FULL ROBUSTNESS TABLE
%% ================================================================
\section{Full Sensitivity Analysis}
\label{app:robustness}

Table~\ref{tab:full_robustness} reports predicted ownership under all robustness specifications.

\begin{table}[htbp]
\centering
\caption{Full Robustness Results}
\label{tab:full_robustness}
\begin{threeparttable}
\begin{tabular}{llc}
\toprule
Category & Specification & Ownership (\%) \\
\midrule
\multicolumn{3}{l}{\textit{Risk aversion ($\gamma$)}} \\
 & $\gamma = 1.50$ & 0.0 \\
 & $\gamma = 2.00$ & 0.0 \\
 & $\gamma = 2.30$ & 0.0 \\
 & $\gamma = 2.40$ & 0.0 \\
 & $\gamma = 2.45$ & 0.0 \\
 & $\gamma = 2.50$ (baseline) & 1.4 \\
 & $\gamma = 2.55$ & 5.3 \\
 & $\gamma = 2.60$ & 6.9 \\
 & $\gamma = 2.70$ & 12.4 \\
 & $\gamma = 2.80$ & 17.1 \\
 & $\gamma = 3.00$ & 23.7 \\
 & $\gamma = 3.50$ & 31.6 \\
 & $\gamma = 4.00$ & 35.3 \\
 & $\gamma = 5.00$ & 39.6 \\[4pt]
\multicolumn{3}{l}{\textit{Hazard multipliers}} \\
 & $[0.45, 1.0, 3.5]$ (R-S functional) & 0.0 \\
 & $[0.50, 1.0, 3.0]$ (baseline) & 1.4 \\
 & $[0.57, 1.0, 2.7]$ (HRS SRH) & 5.7 \\
 & $[0.60, 1.0, 2.0]$ (conservative) & 24.4 \\[4pt]
\multicolumn{3}{l}{\textit{Minimum purchase threshold}} \\
 & \$0 & 1.4 \\
 & \$1,000 & 1.4 \\
 & \$5,000 & 1.4 \\
 & \$10,000 & 1.4 \\
 & \$25,000 & 1.4 \\[4pt]
\multicolumn{3}{l}{\textit{Path dependence}} \\
 & Ordering A & 1.4 \\
 & Ordering B & 1.4 \\[4pt]
\multicolumn{3}{l}{\textit{Grid convergence}} \\
 & Coarse ($40 \times 15$) & 13.6 \\
 & Medium ($60 \times 20$) & 6.5 \\
 & Production ($80 \times 30$) & 1.4 \\
 & Fine ($100 \times 40$) & 1.4 \\[4pt]
\multicolumn{3}{l}{\textit{Bequest specification}} \\
 & No bequests ($\theta = 0$) & 15.7 \\
 & DFJ luxury ($\theta = 56.96$, $\kappa = \$272$K) & 1.4 \\
\bottomrule
\end{tabular}
\begin{tablenotes}
\small
\item Unless otherwise noted, all specifications use baseline parameters: $\gamma = 2.5$, $\beta = 0.97$, DFJ bequests ($\theta = 56.96$, $\kappa = \$272{,}628$), MWR $= 0.82$, $F = \$1{,}000$, $\pi = 2\%$, hazard multipliers $[0.50, 1.0, 3.0]$, production grid ($80 \times 30 \times 101$).
\end{tablenotes}
\end{threeparttable}
\end{table}


%% ================================================================
%%  G. SOLUTION ALGORITHM DETAILS
%% ================================================================
\section{Solution Algorithm}
\label{app:algorithm}

\subsection{Backward induction}

The model is solved by backward induction from the terminal period $T = 46$ (age 110) to $t = 1$ (age 65). At the terminal period, the individual consumes all available resources and leaves residual wealth as a bequest:
\[
V(W, A, H, T) = \max_{c \leq W + A_T + \text{SS}(T) - m_T} \Big\{ u(c) + v\big((1+r)(W + A_T + \text{SS}(T) - m_T - c)\big) \Big\}.
\]

For $t = T-1, \ldots, 1$, the value function is computed at each grid point $(W_i, A_j, H_k)$ by solving the one-period consumption problem using Brent's method on the interval $[\underline{c}, \; W + A_t + \text{SS}(t) - m]$. The continuation value $V(W', A', H', t+1)$ is evaluated by piecewise linear interpolation over the wealth grid with flat extrapolation beyond the grid boundaries.

\subsection{Expectation computation}

The expectation over next-period health and medical expenditures involves:
\begin{enumerate}
\item A sum over three health states weighted by transition probabilities $\pi(H_t, H', t)$.
\item A sum over five Gauss--Hermite quadrature nodes for the medical expenditure shock $\varepsilon$.
\end{enumerate}
Each grid point requires $3 \times 5 = 15$ evaluations of the continuation value, yielding a total of $80 \times 30 \times 3 \times 46 \times 15 \approx 49.7$ million value function evaluations per model solve.

\subsection{Annuitization optimization}

At $t = 1$, the optimal annuitization fraction $\alpha^*$ is found by grid search over 101 equally spaced values in $[0, 1]$. For each candidate $\alpha$, post-purchase wealth is $(1 - \alpha)W_0$ and annuity income is $\alpha W_0 \times \text{payout\_rate}$. The value function at the resulting state is evaluated by bilinear interpolation over the $(W, A)$ grid for the relevant health state. The fixed cost $F$ is subtracted from value whenever $\alpha > 0$.

\subsection{Population evaluation}

Ownership is evaluated on a population of HRS individuals with observed initial wealth, permanent income, and age. For each individual, the optimal $\alpha^*$ is computed at their specific $(W_0, y, H_0)$. Ownership is defined as $\alpha^* > 0$. The ownership rate is the fraction of the eligible population (wealth $\geq$ \$5,000) that annuitizes any positive amount.

\subsection{Computational performance}

A single model solve (one parameterization) requires approximately 60--120 seconds on an Apple M1 processor. The full decomposition (6 solves) runs in 6--12 minutes. The robustness analysis (approximately 40 configurations) requires 1--2 hours. Grid convergence checks (8 configurations at varying resolutions) require 2--4 hours.


\bibliographystyle{aer}
\bibliography{bibliography}

\end{document}
